\section{Canonical bundle data and the least-index root cover}\label{sec:canonical}
\label{altord:sec:canonical}

We return to the original algebraic fibre space. The goal of this
section is to construct the rational line to which
Theorem~\ref{thm:adjoint} applies and to retain the exact section
comparison with the source. The generic fibre of a relative Iitaka
fibration has Kodaira dimension zero; its least-index canonical
root cover provides the entire highest Hodge line.

We work over \(\C\). All varieties in this section are projective.
A rational pair $(Y,B)$ is klt if, on a log resolution, the boundary
defined by crepant pullback of $K_Y+B$ has every coefficient less
than one; it is log canonical if those coefficients are at most
one. We allow negative coefficients when explicitly discussing a
sub-pair. A b-divisor records compatible divisor traces on all
birational models of the base. Saying that its moduli part descends
to $Y$ means that every higher trace is the pullback of its trace $M$
on $Y$.
We distinguish a divisor exceptional over a birational model of the
total space from a divisor whose image on the base has codimension at
least two.

\subsection{The canonical bundle formula and its normalization}

We use the following parabolic form of the canonical bundle formula as
a deep theorem.

\begin{theorem}[Parabolic canonical bundle formula]
\label{altord:cb:parabolic}
Let \(g_0:X_0\to Y_0\) be an algebraic fibre space between smooth
projective varieties, with geometric generic fibre of Kodaira dimension
zero.  After birational modifications of its source and base there is a
commutative diagram
\[
\begin{array}{ccc}
 X'&\xrightarrow{\alpha}&X_0\\
 {\scriptstyle g}\downarrow&&\downarrow{\scriptstyle g_0}\\
 Y&\xrightarrow{\beta}&Y_0
\end{array}
\]
with \(X'\) and \(Y\) smooth, and rational divisors \(B,M,R\) such that
\begin{equation}
\label{altord:cb:formula}
 K_{X'}\sim_{\Q}g^*(K_Y+B+M)+R.
\end{equation}
Here \(B\geq0\), its support has simple normal crossings, \((Y,B)\) is
klt, and \(M\) is nef.  Moreover,
\begin{equation}
\label{altord:cb:residual}
 g_*\mathcal O_{X'}(\lfloor kR^+\rfloor)=\mathcal O_Y
 \quad(k\in\Z_{>0}),
 \qquad
 \operatorname{codim}_Y g(\Supp R^-)\geq2,
\end{equation}
and \(R^-\) is \(\alpha\)-exceptional.
The divisor \(M\) is the trace of the moduli b-divisor, which can be
assumed to descend to \(Y\).  On a sufficiently high generically finite
base change admitting semistable reduction, its pullback is
\(\Q\)-linearly equivalent to a positive rational multiple of the
Schmid extension of the distinguished canonical Hodge eigensheaf of
the canonical root cover.
\end{theorem}

The residual divisor and discriminant assertions, with the normalization
by the least nonzero pluricanonical index, are
\cite[\S4.4 and Theorem~4.5(i)--(v)]{FujinoMori}.
For original boundary zero, the logarithmic semistable part agrees
with the semistable part of Section~2 by
\cite[Proposition~4.7]{FujinoMori}.
The moduli interpretation follows from the discriminant-threshold
definition in \cite[Definition~4.3]{FujinoMori} and the following Ambro
results.  The b-divisor and Hodge assertions use
\cite[Theorems~0.2 and~4.4(1),(3), Lemma~5.2(4)--(5), and
Propositions~5.4--5.5]{AmbroBoundary};
see also \cite[Proposition~3.1]{AmbroModuli}.
These statements allow the auxiliary sub-boundary to have negative
coefficients.

We record how the hypotheses of these results fit together.
In the notation of \cite[\S4.4]{FujinoMori}, applied with original boundary
zero, the crepant boundary upstairs is the negative of the effective
relative canonical divisor over the smooth original source.
Thus the choice \(\Xi=0\) satisfies the hypotheses of their
Theorem~4.5, including those of part~(v).
The coefficients of our \(B\) are their \(s_P/b\); part~(v) places them
in \([0,1)\).  Their residual divisor is \(bR\).
We verify the all-degree direct-image condition from the normalization
below, so that it remains available after the final change of model.
For the exceptionality assertion, flatten before resolving the main
transform, as in \cite[\S4.4]{FujinoMori}.  The inverse image of a base
subset of codimension at least two has codimension at least two in
the flat main transform; finite normalization preserves this bound.
A prime divisor of its resolution lying over that subset is therefore
exceptional over the normalized main transform.  This transform maps
birationally to \(X_0\), so the prime is also \(\alpha\)-exceptional.
Apply this to every component of \(R^-\).  Its small image on \(Y\)
alone would not imply exceptionality over \(X_0\).

Here is also the relevant normalization of the associated sub-pair.
On the geometric generic fibre \(C\), let \(b\) be the least positive
integer with \(H^0(C,bK_C)\ne0\), and let \(D_C\) be the divisor of its
unique section up to scalar.  The horizontal restriction is
\[
 R_C=\frac1bD_C,\qquad
 \Delta_C=-R_C.
\]
The least positive integer making \(K_C+\Delta_C\) principal is
also \(b\). Indeed, \(bK_C-D_C\) is principal. Conversely, if
\(m(K_C-D_C/b)\) is principal for a positive integer \(m\), then
\(mD_C/b\) is an integral effective divisor linearly equivalent to
\(mK_C\). Thus \(H^0(C,mK_C)\ne0\), which forces \(m\geq b\).
This identifies the least-index normalization in the sub-pair
construction with the pluricanonical index used for the root cover.
A rational relative pluricanonical section first defines a divisor
\(R_0\) with this restriction.  For a prime divisor \(P\subset Y\),
subtract from \(R_0\) the pullback of \(P\) multiplied by
\[
 a_P=\min_{E\mapsto P}
       \frac{\operatorname{coeff}_E(R_0)}
            {\operatorname{mult}_E(g^*P)}.
\]
Only finitely many \(a_P\) are nonzero.  The resulting residual divisor
is nonnegative over the generic point of \(P\), with coefficient zero
on at least one component over \(P\).

This normalization also gives the direct-image condition in every
positive degree. Let \(V\subset Y\) be a nonempty open subset, let
\(k>0\) be an integer, and let \(\varphi\) represent a section of
\(\mathcal O_{X'}(\lfloor kR^+\rfloor)\) over \(g^{-1}(V)\).
On the geometric generic fibre \(C\), its poles are bounded by
\(\lfloor kD_C/b\rfloor\), hence by an integral multiple of
\(D_C\). Since \(D_C\sim bK_C\) and \(\kappa(C)=0\), every such
function is constant on \(C\). Since \(g_*\mathcal O_{X'}=\mathcal O_Y\),
the field \(\C(Y)\) is algebraically closed in \(\C(X')\).
Thus \(\varphi=g^*\psi\) for a rational function \(\psi\) on
\(V\). For each prime divisor \(P\) meeting \(V\),
choose a component \(E\) over \(P\) on which the normalized
residual coefficient is zero. Regularity of \(\varphi\) along
\(E\) gives
\[
 0\leq\operatorname{ord}_E(g^*\psi)
   =\operatorname{mult}_E(g^*P)\operatorname{ord}_P(\psi).
\]
Thus \(\psi\) has no codimension-one pole on \(V\), and normality
of \(Y\) makes it regular. Pullbacks of regular functions give the
reverse inclusion, proving
\[
 g_*\mathcal O_{X'}(\lfloor kR^+\rfloor)=\mathcal O_Y
 \qquad(k\in\Z_{>0}).
\]

For the sub-pair \((X',-R)\) its discriminant coefficient is \(1-t_P\),
where
\[
 t_P=\sup\{t\in\R:
       (X',-R+t\,g^*P)\text{ is sub-lc over the generic point of }P\}.
\]
Nonnegativity of \(R\) there implies \(t_P>0\); the component with
residual coefficient zero gives \(t_P\leq1\).
This explains the effective boundary in this normalization.

The rank-one hypothesis for this sub-pair holds even though
\(\Delta_C\) is negative.  Indeed, \(\kappa(C,R_C)=0\).
On a log resolution of \((C,-R_C)\), a section of the rounded
discrepancy sheaf pushes down to a rational function on \(C\) whose
poles are bounded by a sufficiently large integral multiple of
\(D_C\).  The space of such functions has dimension one, while the
constant function is allowed.  Consequently
\[
 \rk g_*\mathcal O_{X'}\bigl(\lceil\mathbf A(X',-R)\rceil\bigr)=1,
\]
where the expression is interpreted by pushing down from a log
resolution.  This is the rank condition in
\cite[introductory assumptions (1)--(3)]{AmbroBoundary}.

After choosing a model on which the moduli b-divisor descends, take
a common higher base model satisfying the flattening and
normal-crossing requirements, with boundary containing the total
transform of the original normalized discriminant support and all
exceptional base divisors. Crepantly pull back the auxiliary sub-pair
to a smooth resolved main transform, taking a log resolution that
includes its residual support and the pullbacks of the base boundary.
Apply the same minimum-coefficient normalization to this residual divisor. Its
horizontal restriction is again \(D_C/b\) for the unique
least-index pluricanonical divisor on the new geometric generic
fibre: for the induced birational morphism \(\lambda:C'\to C\),
crepant pullback gives
\[
 \lambda^*(D_C/b)+K_{C'/C}=D_{C'}/b.
\]
The preceding argument therefore restores the all-degree
direct-image condition on this final model. The normalized residual
is nonnegative over every codimension-one base point; its negative
part has base image of codimension at least two and is exceptional
over the original source by the flattening argument above. Over the
strict transform of an original base prime, its zero-coefficient
component survives and the crepant residual stays nonnegative, so
renormalization makes no further change there. Any new discriminant
support is therefore exceptional and lies in the chosen
simple-normal-crossing boundary. Its coefficients are in \([0,1)\)
by the same threshold calculation, so the discriminant is effective
and klt.

This renormalization does not change the moduli b-divisor. Indeed,
changing a sub-boundary from \(\Delta\) to \(\Delta+g^*A\) changes
the divisor pulled back in \(K_{X'}+\Delta\) by \(A\), while the
discriminant changes by \(A\) as well, by its threshold definition.
Their difference is unchanged, and the same calculation holds on
every higher base model. Crepant birational pullback likewise
preserves the induced b-divisor. Thus the final normalized residual
retains the descended nef moduli part together with the effective
klt discriminant and the residual properties needed for the section
comparisons below.
The parabolic pullback statement and independence of the chosen root
trivialization are also given directly in
\cite[Definition-Proposition~7.1]{AmbroBoundary}.

\begin{remark}
\label{altord:cb:no-abundance}
No abundance assertion for \(M\) is part of the input.
The stronger conclusion of
\cite[Theorem~3.3]{AmbroModuli} assumes effectivity of the boundary on
the geometric generic fibre, an assumption which need not hold for
\(-R_C\).  The b-nefness and Hodge identification used above instead
come from the results of \cite{AmbroBoundary} just cited.
In particular, no good minimal model of \(C\) is required.
For the Hodge-theoretic formulation with sub-pairs, see also
\cite[Definition~6.18, Construction-Definition~6.23,
Remark~6.25, and Theorem~6.28]{BFMT}.
\end{remark}

\subsection{The least-index root cover and its rational variation}

\begin{lemma}[Canonical root cover]
\label{altord:cb:root-cover}
Let \(C\) be a smooth connected projective variety with
\(\kappa(C)=0\), and let \(b>0\) be the least integer for which
\(H^0(C,bK_C)\ne0\).  Choose a nonzero section \(s\) of \(bK_C\) and a
nonzero rational canonical form \(\omega\).
Let \(T\) be the normalization of \(C\) in the field obtained by
adjoining a root
\[
 \xi^b=\frac{s}{\omega^b}.
\]
Then \(T\to C\) has degree \(b\), and every smooth projective
resolution \(W\) of \(T\) satisfies
\[
 \kappa(W)=0,\qquad h^0(W,K_W)=1.
\]
The unique canonical line is spanned by the form
\(\xi h^*\omega\), where \(h:W\to C\) is the induced morphism.
\end{lemma}

\begin{proof}
Put \(n=\dim C\) and \(D_C=\operatorname{div}(s)\).
The case \(n=0\) is immediate.  For \(n>0\), first note that every
nonzero pluricanonical space on \(C\) has dimension one:
two independent sections in the same degree would give a nonconstant
rational map and hence positive Kodaira dimension.

The field \(\C(C)\) contains all roots of unity.
If \(s/\omega^b=u^p\) for a prime \(p\mid b\), the rational
\((b/p)\)-canonical form \(u\omega^{b/p}\) has \(p\)-th power \(s\).
At every prime divisor its order is \(p^{-1}\) times the nonnegative
order of \(s\), so this form has no poles.  A rational section of a
line bundle on the smooth variety \(C\) with no divisorial poles is
regular.  This contradicts the minimality of \(b\).
The Kummer irreducibility criterion now gives a field extension of
degree \(b\).  Denote the normalization map by \(\pi:T\to C\), and a
smooth projective resolution by \(\rho:W\to T\), so \(h=\pi\rho\).

The differential of \(h\) gives a nonzero morphism
\(h^*K_C\to K_W\) at the generic point, and a regular morphism of
line bundles everywhere.  Use this morphism and its tensor powers
when pulling back canonical forms.  The rational canonical form
\(\eta=\xi h^*\omega\) then satisfies
\[
 \eta^{\otimes b}=h^*s.
\]
The right side is a regular section of \(bK_W\).
Thus \(b\,\operatorname{ord}_E(\eta)\geq0\) for every prime divisor
\(E\subset W\).  Normality of \(W\) again implies that \(\eta\) is
regular, and therefore \(h^0(W,K_W)\geq1\).

We next bound the canonical divisor from above.
Choose \(K_C=\operatorname{div}(\omega)\), and on \(T\) use the
canonical Weil divisor determined by the differential pullback of
\(\omega\).  At a prime divisor \(P\subset C\) outside \(\Supp D_C\),
the order of \(s/\omega^b\) is divisible by \(b\).  Locally over the
discrete valuation ring at \(P\), after dividing the root by a
power of a uniformizer, the defining equation is a \(b\)-th root of
a unit.  It is unramified, since the characteristic is zero.
Consequently all ramification divisors lie over \(\Supp D_C\).

If \(Q\subset T\) lies over \(P\subset\Supp D_C\), write
\(e=e(Q/P)\) and \(d=\operatorname{ord}_P(D_C)\geq1\).
Tameness gives the ramification coefficient \(e-1\), and
\[
 e-1\leq ed=\operatorname{ord}_Q(\pi^*D_C).
\]
The canonical divisor formula for finite maps, interpreted at
codimension-one points, therefore gives
\[
 K_T\leq\pi^*(K_C+D_C).
\]
On \(W\) the same inequality holds away from the
\(\rho\)-exceptional divisors.  Hence there is an effective integral
\(\rho\)-exceptional divisor \(E\) such that
\begin{equation}
\label{altord:cb:cover-upper-bound}
 K_W\leq h^*(K_C+D_C)+E.
\end{equation}
This argument does not require a pullback of the Weil divisor \(K_T\).

Set \(A=K_C+D_C\).  Normality of \(T\) implies
\(\rho_*\mathcal O_W(mE)=\mathcal O_T\) for every \(m\geq0\):
a rational function allowed poles only on exceptional divisors is
regular in codimension one on \(T\), hence regular on \(T\).
Projection formula and \eqref{altord:cb:cover-upper-bound} imply
\[
 \kappa(W,K_W)
 \leq \kappa(W,h^*A+E)
 =\kappa(T,\pi^*A)
 =\kappa(C,A).
\]
The last equality is invariance of Iitaka dimension under a finite
surjective pullback.  In this situation it can be seen directly
from section rings: a homogeneous section upstairs satisfies its
monic characteristic polynomial over \(\C(C)\); its coefficients,
in degree \(jm\), are sections of \(jmA\), by the valuation bounds
for sections of \(m\pi^*A\).  Thus the upstairs graded ring is
integral over the downstairs ring, and their fraction fields have
the same transcendence degree.
Finally \(D_C\sim bK_C\), so
\[
 \kappa(C,A)=\kappa(C,(b+1)K_C)=0.
\]
Together with the nonzero form \(\eta\), this proves
\(\kappa(W)=0\), and then \(h^0(W,K_W)=1\).
\end{proof}

The underlying cover of Lemma~\ref{altord:cb:root-cover} is the
canonical root cover used in
\cite[Remark~2.6 and Theorem~3.1]{FujinoMori}.
Over a dense open subset \(U\subset Y\), a rational relative
pluricanonical section constructs this cover in families; a
projective resolution, followed by shrinking \(U\), gives a smooth
projective morphism \(h_U:\mathcal W_U\to U\).
Its geometric generic fibre has the invariants in
Lemma~\ref{altord:cb:root-cover}. Constancy of Hodge numbers in the smooth
projective family then gives geometric genus one on every fibre.
If their dimension is \(n\), then
\[
 \mathbb V=R^n(h_U)_*\Q,\qquad
 \rk F^n(\mathbb V\otimes\mathcal O_U)=1.
\]
The variation has an integral lattice modulo torsion and is
polarizable, using a relative ample class and the Lefschetz
decomposition.  Alternatively, its primitive middle cohomology
is polarizable and contains all of \(H^{n,0}\), since cup product
with a \((1,1)\)-class annihilates \(H^{n,0}\).
The cyclic deck transformation \(\xi\mapsto\zeta\xi\) acts on the
unique top form through \(\zeta\).  Thus the distinguished
canonical eigensheaf is the entire top Hodge line of a rational
polarizable variation, even when the character itself is not
defined over \(\Q\).  We do not replace the rational variation by
that character eigenspace.

After a smooth projective alteration \(\tau:T_Y\to Y\), resolution
of the boundary, and semistable reduction in codimension one, the
identification in Theorem~\ref{altord:cb:parabolic} reads
\begin{equation}
\label{altord:cb:hodge-realization}
 \tau^*M\sim_{\Q}aJ,\qquad a\in\Q_{>0},
\end{equation}
where \(J\) is the Schmid extension of the top Hodge line of
\(\tau^*\mathbb V\), with unipotent local monodromy.
Birational changes of the resolved root cover do not change its
canonical direct image: on a common resolution, differential
pullback identifies its canonical forms.
The extension comparison is the one in
\cite[Theorem~4.4(1),(3), Lemma~5.2(5), and
Proposition~5.5]{AmbroBoundary}, rather than
an identification asserted only on \(U\).
Changing the chosen rational trivialization by a function from
the base gives the same identification after a further finite
base change extracting the corresponding root.

\subsection{Relative Iitaka reduction and section comparisons}

\begin{proposition}[Canonical bundle reduction]
\label{altord:cb:reduction}
Let \(f:X\to Z\) be an algebraic fibre space between smooth connected
projective varieties, and let its geometric generic fibre \(F\) have
\(d=\kappa(F)\geq0\).
There exist smooth projective varieties \(\widetilde X,Y\),
a birational morphism \(\mu:\widetilde X\to X\), and algebraic fibre
spaces
\[
 \widetilde X\xrightarrow{g}Y\xrightarrow{q}Z,\qquad
 qg=f\mu,
\]
such that \(\dim Y-\dim Z=d\) and the geometric generic fibre of
\(g\) has Kodaira dimension zero.
They admit the data \(B,M,R\) of Theorem~\ref{altord:cb:parabolic},
with \(R^-\) exceptional over \(X\).
The moduli divisor \(M\) has the realization
\eqref{altord:cb:hodge-realization} for a polarizable rational variation
whose top Hodge piece has rank one.
Writing
\[
 D=K_Y+B+M,
\]
one has
\begin{equation}
\label{altord:cb:two-comparisons}
 \kappa(X)\geq\kappa(Y,D),
 \qquad
 D|_{Y_{\bar\eta}}\text{ is big},
\end{equation}
where \(\bar\eta\) is a geometric generic point of \(Z\).
The same bigness holds on very general smooth fibres of \(q\).
\end{proposition}

\begin{proof}
Over \(\C(Z)\), take a sufficiently divisible pluricanonical map
of the generic fibre and its Iitaka fibration, including the
relative algebraic closure of the image field.
The Iitaka fibration theorem gives image dimension \(d\) and
Kodaira dimension zero for the geometric generic fibre of this
fibration.  Formation of its pluricanonical spaces commutes with
extension to an algebraic closure of \(\C(Z)\); thus these
statements can be checked geometrically.
Extend the rational map to projective models over \(Z\), resolve
source and target, and take the connected-fibre factorization.
This gives \(X_0\to Y_0\to Z\) with the required dimensions.
The second map has connected fibres as well: pushing
\(\mathcal O_{X_0}\) first to \(Y_0\) and then to \(Z\) gives
\(\mathcal O_Z\), because the original \(f\) and the birational
map to its smooth source have this property.

Apply Theorem~\ref{altord:cb:parabolic} to \(X_0\to Y_0\).
All its modifications can be performed over \(Z\).
Exceptionality over \(X_0\) implies exceptionality over \(X\).
Lemma~\ref{altord:cb:root-cover} and the subsequent family construction
give the required rational variation and its top line.
The two comparisons are proved in Lemma~\ref{altord:cb:comparison}
below.
\end{proof}

\begin{lemma}[Comparison of section growth]
\label{altord:cb:comparison}
For the factorization and canonical bundle data in
Proposition~\ref{altord:cb:reduction}, \eqref{altord:cb:two-comparisons} holds.
\end{lemma}

\begin{proof}
The formula gives
\[
 K_{\widetilde X}+R^-
 \sim_{\Q}g^*D+R^+.
\]
Write \(K_{\widetilde X}=\mu^*K_X+E_\mu\), using compatible
canonical divisors.  Since \(X\) and \(\widetilde X\) are smooth,
\(E_\mu\) is effective and \(\mu\)-exceptional.
The same is true of \(E_\mu+R^-\).
For sufficiently divisible \(m\), normality and projection
formula therefore give
\[
 H^0\bigl(\widetilde X,m(K_{\widetilde X}+R^-)\bigr)
 \simeq H^0(X,mK_X).
\]
Multiplication by the canonical section of \(mR^+\) injects
\(H^0(Y,mD)\) into the space on the left; here we also use
\(g_*\mathcal O_{\widetilde X}=\mathcal O_Y\).
This proves \(\kappa(Y,D)\leq\kappa(X)\).

Now make the flat base change \(\bar\eta\to Z\).
The varieties \(\widetilde X_{\bar\eta}\) and \(Y_{\bar\eta}\)
are smooth; the former is birational to the geometric generic
fibre \(F\) of \(f\).
Restrictions of canonical divisors agree, up to linear
equivalence, with the corresponding canonical divisors of these
fibres.  The restricted formula consequently implies
\[
 K_{\widetilde X_{\bar\eta}}
 \leq_{\Q}
 g_{\bar\eta}^*(D|_{Y_{\bar\eta}})
 +R^+|_{\widetilde X_{\bar\eta}}.
\]
Here \(\leq_{\Q}\) means that the difference is
\(\Q\)-linearly equivalent to an effective divisor.
Formation of a proper direct image commutes with this flat
base change, so \eqref{altord:cb:residual} yields, for sufficiently
divisible \(m\),
\[
 (g_{\bar\eta})_*
 \mathcal O_{\widetilde X_{\bar\eta}}
       (mR^+|_{\widetilde X_{\bar\eta}})
 =\mathcal O_{Y_{\bar\eta}}.
\]
Projection formula now gives
\begin{align*}
 h^0(\widetilde X_{\bar\eta},mK_{\widetilde X_{\bar\eta}})
 &\leq
 h^0\bigl(\widetilde X_{\bar\eta},
 m g_{\bar\eta}^*(D|_{Y_{\bar\eta}})
       +mR^+|_{\widetilde X_{\bar\eta}}\bigr)\\
 &=h^0(Y_{\bar\eta},mD|_{Y_{\bar\eta}}).
\end{align*}
It follows that
\[
 d=\kappa(F)
 \leq\kappa(Y_{\bar\eta},D|_{Y_{\bar\eta}})
 \leq\dim Y_{\bar\eta}=d.
\]
Thus \(D|_{Y_{\bar\eta}}\) is big, including the
zero-dimensional case.

For completeness, the bigness assertion spreads to a dense open
subset after restricting to the smooth locus of \(q\).
Fix a relatively ample divisor \(H\).
Bigness on the geometric generic fibre gives, for some
sufficiently divisible integer \(m\),
a nonzero section of \(mD-H\) there.  Flat scalar extension of
cohomology first gives such a section over \(\C(Z)\).
Shrinking the base, the section extends to the family and remains
nonzero on its fibres.  Hence \(mD|_{Y_z}\) is the sum of the
ample divisor \(H|_{Y_z}\) and an effective divisor for general
\(z\), proving the stated very-general-fibre assertion.
\end{proof}

\begin{lemma}[Further birational modifications]
\label{altord:cb:high-model}
Suppose \(Y,B,M,D\) satisfy the conclusions of
Proposition~\ref{altord:cb:reduction}.  Let \(\nu:Y'\to Y\) be a smooth
projective log resolution, and write
\[
 K_{Y'}+B^{\mathrm{cr}}=\nu^*(K_Y+B).
\]
Set \(B'=(B^{\mathrm{cr}})^+\), \(M'=\nu^*M\), and
\(D'=K_{Y'}+B'+M'\).
Then \(B'\) is effective, \((Y',B')\) is klt, \(M'\) is nef with
the same Hodge realization after pullback, and
\[
 D'=\nu^*D+(B^{\mathrm{cr}})^-.
\]
The additional divisor is effective and \(\nu\)-exceptional.
Consequently \(\kappa(Y',D')=\kappa(Y,D)\), and \(D'\) is big
on the geometric generic fibre over \(Z\).
\end{lemma}

\begin{proof}
The nonexceptional coefficients of \(B^{\mathrm{cr}}\) are those
of the effective divisor \(B\), so its negative part is
exceptional.  All coefficients of \(B^{\mathrm{cr}}\) are less
than one by the klt hypothesis.  Its positive part has simple
normal-crossing support and coefficients in \([0,1)\), proving
that \((Y',B')\) is klt.
The displayed identity follows directly from the definitions.
Exceptional effective twists do not change section spaces of a
pullback to a normal base, so they do not change the Iitaka
dimension.
On the geometric generic fibre over \(Z\), the induced map is
birational and the restricted identity expresses \(D'\) as the
pullback of a big divisor plus an effective divisor.
Nefness is preserved by pullback.  For the Hodge assertion, use
a common smooth alteration of \(Y'\) and the alteration in
\eqref{altord:cb:hodge-realization}; the unipotent Deligne and Schmid
extensions commute with the resulting morphisms of smooth
normal-crossing pairs.
\end{proof}

\begin{remark}\label{altord:cb:reduced-moduli-remark}
After Lemma~\ref{altord:cb:high-model}, \(B'\) need not be described as
the discriminant of the original sub-pair.  The properties used
later are precisely effectiveness and klt singularities, the
nef Hodge realization of \(M'\), and the two comparisons in
\eqref{altord:cb:two-comparisons}; all have been preserved.
In particular, since the big cone of \(Y_{\bar\eta}\) is open,
there is a rational number \(v\) with \(0<v<1\) such that
\[
 (K_Y+B+vM)|_{Y_{\bar\eta}}
 \quad\text{is big}.
\]
This remains true on very general fibres by the spreading
argument in Lemma~\ref{altord:cb:comparison}.
\end{remark}
